% Chapter 4: Results & Discussion

\chapter{Results \& Discussion}
\noindent
This chapter evaluates the performance of the proposed memetic framework, transitioning from environmental validation and algorithmic calibration on synthetic topologies to the final network optimization of the empirical Iligan City \texttt{TravelGraph}.

% ────────────────────────────────────────────────────────────
\section{Environment and Demand Generation}
\noindent
To ensure the optimization relies on valid spatial and demographic data, the digital twin of Iligan City and the simulated passenger demand must first be verified.

\subsection{\texttt{CityGraph} Arterial Network Pruning}
\textbf{Visual 1:} Full Iligan City OpenStreetMap graph (all \texttt{DirEdge} connections). \\
\textbf{Visual 2:} The filtered Arterial \texttt{CityGraph}. \\
\textit{Discussion:} Proves the successful removal of residential dead-ends and non-traversable nodes, restricting \texttt{Jeep} agent movement to valid, navigable transit corridors while leaving the full network available for passenger walking access.

\subsection{\texttt{DirectDemandModel} (DDM) Topography}
\textbf{Visual 1:} Centrality-weighted nodes queried from the TomTom Traffic Flow API under Efraimidis \& Spirakis weighted-reservoir sampling. \\
\textbf{Visual 2:} The Inverse Distance Weighting (IDW) heatmap showing spatial traffic-weight diffusion. \\
\textbf{Visual 3:} Betweenness Centrality graph vs. the fused OD Centrality graph ($P_i = W_i^{\alpha} C_i^{\beta}$), side by side. \\
\textbf{Visual 4:} Temporal DDM surfaces and their demand distributions for the 08:00, 13:00, and 17:00 query windows. \\
\textbf{Visual 5:} DDM surface projected onto the full network vs. the arterial subset used for waypoint sampling. \\
\textit{Discussion:} Demonstrates that OD Centrality, derived from live traffic fused with network topology, is a superior proxy for passenger flow than raw topological betweenness: corridors that are both structurally central and empirically congested dominate the surface, producing a spatially coherent demand gradient rather than scattered hot spots.

\subsection{Demand-Responsive Route Generation}
\textbf{Visual 1:} Representative generation-zero route loops produced by the \texttt{RouteGenerator}. \\
\textit{Discussion:} Confirms that waypoints sampled from the DDM surface under the \texttt{only\_drivable} constraint yield structurally valid closed loops (passing loop continuity, Layer-2 isolation, and single out-pointer checks), biasing initial routes toward high-demand corridors.

% ────────────────────────────────────────────────────────────
\section{Architectural Validation}
\noindent
Prior to evolutionary optimization, the foundational mechanics of the multi-layer graph and the temporal physics engine are validated.

\subsection{The Three-Layer \texttt{TravelGraph} Connectivity}
\textbf{Visual 1:} A matrix of six isolated illustrations detailing the explicit layer transitions: Start Walk (SW), Wait (WA), Ride (RI), Alight (AL), End Walk (EW), and Transfer (TR). \\
\textbf{Visual 2:} A sample \texttt{Journey} tracing a single \texttt{Passenger} agent navigating physical distance and transfer penalties, converted into Equivalent In-Vehicle Minutes (EIVM). \\
\textit{Discussion:} Demonstrates that the A* pathfinder optimizes generalized cost across a single EIVM scale rather than raw geometric distance, and that route isolation prevents the ``teleportation anomaly'' between physically coincident routes.

\subsection{Agent-Based Temporal Simulation}
\textbf{Visual 1:} Tick-by-tick snapshots tracing \texttt{Passenger} and \texttt{Jeep} agents, showing boarding, capacity-limited waiting, alighting, and headway behavior. \\
\textit{Discussion:} Validates that the engine reproduces the compounding micro-level effects that justify an agent-based model over static assignment: a full jeepney leaves passengers stranded in the \texttt{WAITING} state (accruing delay), while equidistant spawning suppresses early bus bunching.

% ────────────────────────────────────────────────────────────
\section{Component Calibration and Verification}
\noindent
Prior to executing the framework on the full empirical data, hyperparameter calibration was conducted on a synthetic Manhattan grid topology to isolate algorithmic behavior from geographic anomalies \citep{kepaptsoglou2009transit}.

\subsection{Mohring Fleet Allocation Stability}
\textbf{Visual:} Coefficient of variation (CV) of the per-route Mohring fleet allocation as a function of the number of routes in the system, with a horizontal acceptance threshold drawn at $\mathrm{CV} = 0.5$. \\
\textit{Discussion:} The square-root allocation must remain stable regardless of how many routes a candidate system contains. Sweeping the route count and measuring the CV of the allocation across independent demand samples, every tested system size remains below the $0.5$ acceptance threshold, confirming that fleet distribution is robust to system size at the production sample size and does not swing with the stochastic demand draw.

\subsection{Temporal Discretization ($\Delta t$)}
\textit{Discussion:} All production runs fix the simulation time step at $\Delta t = 10\text{s}$ (\texttt{seconds\_per\_tick}). The only discretization error arises from sub-tick rounding when a vehicle reaches a stop; this error is unbiased, so across the large stochastic passenger population and the extended simulation horizon, individual overshoot and undershoot are independent and average toward zero. The residual noise is therefore dominated by, and absorbed into, the stochastic demand variance the optimizer already tolerates, while a 10-second step keeps per-evaluation runtime tractable. A standalone discretization sweep is consequently unnecessary.

\subsection{Simulation Horizon and Demand-Volume Calibration}
\textbf{Visual 1:} Passenger completion rate as a function of the simulation horizon \texttt{num\_ticks}, recorded cumulatively within a single long run. \\
\textbf{Visual 2:} A line chart with shaded error bands ($\pm 1$ standard deviation) plotting passenger volume against mean simulation fitness ($F_{\text{sim}}$) and fitness variance ($\sigma^2$) across repeated runs of a static topology. \\
\textit{Discussion:} Two coupled knobs govern both fidelity and per-evaluation runtime. The horizon must be long enough for the large majority of passengers to complete their journeys, since incomplete passengers are scored through the underservice penalty rather than realized cost; the production horizon is taken at the tick where the completion curve plateaus. The spawn rate is calibrated against fitness variance: under low passenger volumes the score is dominated by stochastic OD-sampling noise, producing extreme variance between identical runs, whereas the variance collapses past a critical volume. The production rate is selected just beyond this collapse, guaranteeing the \texttt{MemeticEngine} selects parents on true structural fitness rather than sampling luck while bounding the $\mathcal{O}(N)$ per-tick cost.

\subsection{Opportunistic Riding and Individual Travel-Delay Benefit}
\textbf{Visual:} Distribution of the travel-time delta $\Delta T = T_{\text{actual}} - T_{\text{expected}}$ for passengers who waited for their expected jeepney (Group~A) versus those who boarded an earlier alternative (Group~B), annotated with the Mann--Whitney $U$ test result. \\
\textit{Discussion:} A positive boarding \texttt{weight\_tolerance} permits a passenger to board an acceptable earlier-arriving jeepney instead of strictly waiting for the single vehicle on their pre-computed least-cost journey. A simulation run with \texttt{weight\_tolerance} $> 0$ partitions completed passengers into Group~A (waited) and Group~B (rode opportunistically), and their travel-time deltas are compared with the non-parametric Mann--Whitney $U$ test. A statistically significant result with a lower median delay for Group~B proves that opportunistic riders experience a real reduction in travel delay, validating that the weight-tolerance mechanism confers an individual-level benefit rather than merely perturbing path choice. The underlying travel-graph cost weights were externally validated and are not re-tuned here.

\subsection{Surrogate Ranking Fidelity (Spearman $\rho$)}
\textbf{Visual:} Scatter plot mapping the $\mathcal{O}(1)$ static A* surrogate cost ($C_{\text{sur}}$) against the full temporal simulation fitness ($F_{\text{sim}}$). \\
\textit{Discussion:} The surrogate is used only as an acceptance gate during Lamarckian local search, so its legitimacy rests on preserving the relative \emph{rank} of candidate solutions rather than their absolute cost. The Spearman ($\rho$) and Kendall ($\tau$) rank correlations, together with top-tier selection precision and recall, confirm that the surrogate never discards a truly optimal candidate, validating its use as a lightweight computational bypass.

\subsection{Lamarckian Operator Mechanics}
\textbf{Visual 1:} Pre-mutation vs. post-mutation snapshots on the Manhattan toy network demonstrating the three local-search mutations: Attraction toward high-demand corridors, Repulsion from oversupplied corridors, and Tortuosity reduction (pruning). \\
\textbf{Visual 2:} Scatter plot proving that networks with lower Demand-Service Gaps natively produce better $F_{\text{sim}}$ scores. \\
\textit{Discussion:} Isolates the structural effect of each operator and confirms that the demand-service gap is the correct signal for the operators to act on.

% ────────────────────────────────────────────────────────────
\section{Evolutionary Dynamics}
\noindent
This section evaluates the behavioral dynamics of the \texttt{MemeticEngine} over time.

\subsection{Convergence and \texttt{AdaptiveController} Orchestration}
\textbf{Visual:} The convergence curve tracking best, average, and worst fitness scores, overlaid with the dynamically scaling mutation rate triggered by the stagnation counter. \\
\textit{Discussion:} Demonstrates that the adaptive controller escalates exploration when the population stagnates and relaxes it as fitness improves, driving total user cost down across generations.

\subsection{Topological Hub Crossover and Lineage Tracking}
\textbf{Visual:} The topological hub and a generational lineage tree from the \texttt{TelemetryEngine}, showing the preservation of high-efficiency ``trunk'' route segments across parent--child generations. \\
\textit{Discussion:} Validates the Topological Hub Crossover by proving that fitter parents disproportionately pass down their core, high-efficiency trunk routes.

\subsection{Epigenetic Pheromone Inheritance}
\textbf{Visual 1:} Stigmergic pheromone surfaces for two parents and the fitness-weighted blended matrix ($w_A \tau^A + w_B \tau^B$) inherited by their offspring. \\
\textbf{Visual 2:} A dual-axis comparison of post-mutation $F_{\text{sim}}$ improvement and computational runtime for offspring mutated with the inherited matrix (Twin~A) versus offspring forced to run a full native pre-simulation (Twin~B). \\
\textit{Discussion:} An A/B cloning experiment shows that Epigenetic Inheritance captures the vast majority of the fitness improvement achieved by a native pre-simulation but executes in a fraction of the time, proving that transferring the cultural demand memory is a Pareto-efficient and computationally mandatory heuristic bypass.

\subsection{Multi-Dimensional Phenotypic Convergence}
\textbf{Visual:} A multi-line plot tracking the population's Jaccard Similarity, Graph Edit Distance (GED), and 2D Wasserstein Distance, overlaid with collapsing fitness variance ($\sigma^2(F)$). \\
\textit{Discussion:} Proves the algorithm avoids flat fitness landscapes by demonstrating simultaneous convergence across three orthogonal dimensions: shared streets (Jaccard), topological connectivity (GED), and spatial demand coverage (Wasserstein).

% ────────────────────────────────────────────────────────────
\section{Optimized Network for Iligan City}
\noindent
With the framework validated, the \texttt{Optimizer} is applied to the empirical Iligan City data. Because no digitized record of the current jeepney routes exists, a stochastic baseline of generation-zero networks establishes the reference for arbitrary-network performance.

\subsection{Stochastic Baseline Generation}
\textit{Discussion:} Evaluates the average performance of randomly initialized generation-zero networks to establish a statistical baseline against which algorithmic improvement is measured.

\subsection{The Optimized Network and Cross-Run Robustness}
\textbf{Visual 1:} A map of a representative optimized route system, reconstructed from the \texttt{StatePreservationEngine} JSON snapshots. \\
\textbf{Visual 2:} A pairwise similarity matrix comparing the final networks from the seven identical-configuration runs (profile P1, varying only the random seed) using Jaccard, GED, and 2D Wasserstein metrics. \\
\textit{Discussion:} Because the optimizer is a stochastic metaheuristic, profile P1 is executed seven times to establish reproducibility. High structural similarity across all seven runs proves that the algorithm converges on a robust optimum region rather than a single lucky seed, and the reduction in total user cost relative to the stochastic baseline quantifies the optimization gain.

\subsection{Robustness to Demand Regime}
\textbf{Visual:} Similarity matrix comparing the network re-optimized under the alternative temporal demand surfaces (08:00, 13:00, 17:00) against the primary optimum, using the Jaccard, GED, and Wasserstein metrics. \\
\textit{Discussion:} Tests whether the optimum is specific to the demand regime it was trained on. High cross-regime similarity indicates a temporally robust trunk structure that serves peak and off-peak demand alike; low similarity indicates that distinct time-of-day networks are warranted. The alternative demand surfaces reuse the temporal DDM pickles, so this axis requires only re-optimization, not new data collection.

% Optional (include only if compute time allows; otherwise defer to Future Work):
% \subsection{Operational Sensitivity}
% \textit{Discussion:} Observes structural changes in the optimized topologies when varying
% operational constraints, namely the allocatable \texttt{Jeep} fleet size and passenger
% spawning volume. This axis overlaps with the preliminary calibration work and is the
% lowest-novelty sensitivity dimension.

\subsection{Equity and Distribution Analysis}
\textbf{Visual:} Histogram comparing the distribution of passenger travel times ($T_i$) between the stochastic baseline and the final optimized network. \\
\textit{Discussion:} Proves the efficacy of the Equity Regularizer ($\alpha \sigma$) by showing a reduction in the long tail of extreme commute times for peripheral communities. This analysis is computed from the seven P1 runs and requires no additional optimization.

\subsection{Network Resilience and Path Diversity}
\textbf{Visual:} A heatmap of edge traversals accompanied by the Shannon Entropy of the passenger path distribution. \\
\textit{Discussion:} Demonstrates that the optimized network avoids forcing all traffic through a single catastrophic bottleneck. A higher Shannon entropy indicates distinct routing options and a resilient transit architecture. Like the equity analysis, this is derived from the existing P1 runs.
